% ---- pg-528 （书页 516） ----

\noindent 11.30\quad 格里尔德和兰茨(见文献[13])对三个砂岩区的原油样本作了化学分析. 这三个地区分别记为 $\pi_1$, $\pi_2$ 和 $\pi_3$. 化学分析的结果是得到样本中的三种痕量元素及两种烃类化合物的含量:

\[
\begin{gathered}
X_1=\text{钒(灰分中百分比)}\\
X_2=\text{铁(灰分中百分比)}\\
X_3=\text{镍(灰分中百分比)}\\
X_4=\text{饱和烃(面积百分比)}\\
X_5=\text{芳香烃(面积百分比)}
\end{gathered}
\]

得到的数据列于\hyperref[tab:11.7]{表 11.7}.

(a) 假定有正态性, 试得出估计的最小 TPM 法则, 对正态性假定的适当性作评述.

(b) 用拉亨布鲁克的提留方法确定 $E(\mathrm{AER})$ 的估计值. 同时给出模糊矩阵.

(c) 考虑能将数据转化到具有正态性的各种变换(见例 11.14), 重复(a)和(b).

\begin{mvtable}
\phantomsection\label{tab:11.7}\mvtabcap{表 11.7\quad 原油数据}
\begin{tabular}{lccccc}
\toprule
& $x_1$ & $x_2$ & $x_3$ & $x_4$ & $x_5$\\
\midrule
$\pi_1$ & 3.9 & 51.0 & 0.20 & 7.06 & 12.19\\
& 2.7 & 49.0 & 0.07 & 7.14 & 12.23\\
& 2.8 & 36.0 & 0.30 & 7.00 & 11.30\\
& 3.1 & 45.0 & 0.08 & 7.20 & 13.01\\
& 3.5 & 46.0 & 0.10 & 7.81 & 12.63\\
& 3.9 & 43.0 & 0.07 & 6.25 & 10.42\\
& 2.7 & 35.0 & 0.00 & 5.11 & 9.00\\
$\pi_2$ & 5.0 & 47.0 & 0.07 & 7.06 & 6.10\\
& 3.4 & 32.0 & 0.20 & 5.82 & 4.69\\
& 1.2 & 12.0 & 0.00 & 5.54 & 3.15\\
& 8.4 & 17.0 & 0.07 & 6.31 & 4.55\\
& 4.2 & 36.0 & 0.50 & 9.25 & 4.95\\
& 4.2 & 35.0 & 0.50 & 5.69 & 2.22\\
& 3.9 & 41.0 & 0.10 & 5.63 & 2.94\\
& 3.9 & 36.0 & 0.07 & 6.19 & 2.27\\
& 7.3 & 32.0 & 0.30 & 8.02 & 12.92\\
& 4.4 & 46.0 & 0.07 & 7.54 & 5.76\\
& 3.0 & 30.0 & 0.00 & 5.12 & 10.77\\
$\pi_3$ & 6.3 & 13.0 & 0.50 & 4.24 & 8.27\\
& 1.7 & 5.6 & 1.00 & 5.69 & 4.64\\
& 7.3 & 24.0 & 0.00 & 4.34 & 2.99\\
& 7.8 & 18.0 & 0.50 & 3.92 & 6.09\\
& 7.8 & 25.0 & 0.70 & 5.39 & 6.20\\
& 7.8 & 26.0 & 1.00 & 5.02 & 2.50\\
& 9.5 & 17.0 & 0.05 & 3.52 & 5.71\\
& 7.7 & 14.0 & 0.30 & 4.65 & 8.63\\
& 11.0 & 20.0 & 0.50 & 4.27 & 8.40\\
\bottomrule
\end{tabular}
\end{mvtable}
